\magicSection{Polynomial Specification}{sec:PolynomialSpec}

We saw in Section~\ref{sec:Tensors} that tensor strides imply a polynomial function parameterized over free variables and an opaque function \texttt{DivAndMod(x, D, M)} $= \frac{x}{D} \% M$.
We can stick with the polynomial theme for specifying the high-level algorithm.

To have polynomials, we only need to have the data type satisfy semi-ring axioms (of course floats actually don't satisfy this, but it's customary to ignore this), with a
\begin{itemize}
  \item ring addition operator
  \item ring multiplication operator
\end{itemize}
Commonly, these are regular addition and multiplication, but the tropical semi-ring is also frequently used, with minimum or maximum being the ring addition operator, and regular addition being the ring multiplication operator.

An \myKeyA{extended polynomial function} is defined to use some semi-ring, and is one of the cases
\begin{itemize}
  \item A constant integer or variable ($N = \sum_{i=1}^N 1$ where $1$ is the multiplicative identity)
  \item Sum of two extended polynomial functions using the same semi-ring
  \item Product of two extended polynomial functions using the same semi-ring
  \item Finite sum $\sum_{k=0}^{B(...)} A(...)$ where $k$ is a new variable and $A$ and $B$ are extended polynomial functions whose parameters are also extended polynomial functions, all using the same semi-ring.
  \item An \myKeyA{opaque function} whose parameters are extended polynomial functions.
    The opaque function may be used to represent indexing and reading an input tensor, transcendental functions, or extended polynomial functions that need not use the same semi-ring.
  \item Reciprocal of such an opaque function (useful for online softmax adjustment)
\end{itemize}
The opaque function and reciprocal are what makes this ``extended''.
Without them, the polynomial function equality problem is decidable, by using Faulhaber's formula for the finite sum, and expanding recursive polynomial functions.
For the prototype, a lot of symbolic evaluation will come down to checking equality of these extended polynomial functions, which represent either addresses or data values.

I am not sure if the added opaque function makes this undecidable.
But pragmatically, we can approximate the equality problem by generating new free variables for each unique opaque function and substituting a non-linear polynomial for each opaque function, e.g.
\begin{equation*}
    A(x_1, x_2, x_3, ...) = \sum_{i} A_i (x_i) (x_i + 1) \text{  where $A_i$ are the generated free variables}
\end{equation*}
For algorithms without data-dependent indexing or branching, it should be possible to specify the expected output buffer values an extended polynomial function of the input buffers.
For example, dual-gemm $C = A(B_1 + B_2)$ would be specified as
\begin{equation*}
    C(m, n) = \sum_{k=0}^{K-1} A(m, k)(B_1(k, n) B_2(k, n))
\end{equation*}
where $A$, $B_1$, $B_2$ are opaque functions.

As mentioned, I'd like to keep an eye out for the future and consider irregular data access.
This polynomial stuff is just a seed idea to get started on thinking.
I don't know if this is really a good approach.
